% GENERATED FILE. Edit canonical lessons, then run npm run generate:books.
% canonical-source-hash: fnv1a64:78bbb599d96e8557
% canonical-lessons: HI-C291-casmdid, HI-C291-bayan, HI-C291-pahcanna, HI-C291-sak, HI-C291-ghurna

\chapter{What the Witness Saw}
\label{ch:hi-what-the-witness-saw}
\hlchaptermodality{291}

\begin{chapteropening}
\textbf{By the end of this chapter:} \emph{I can say what I saw and what would have been different.}

Complete the last lesson of chapter 291: I can say what I saw and what would have been different.
\end{chapteropening}

\section[\texorpdfstring{\hi{चश्मदीद} (caśmdīd)}{चश्मदीद (caśmdīd)}]{\hi{चश्मदीद} (caśmdīd) — an eyewitness}

\label{lesson:HI-C291-casmdid}



\begin{quote}
Before the new one: say the Hindi for a police inspector, then the Hindi for a lock-up.
\end{quote}

\subsection*{You'll want to know: \hi{चश्मदीद}}
\textbf{\hi{चश्मदीद}} — \emph{caśmdīd} — \textquotedblleft{}an eyewitness\textquotedblright{}.

\begin{quote}
\hi{अगर} \hi{कोई} \hi{चश्मदीद} \hi{होता} \hi{तो} \hi{पुलिस} \hi{का} \hi{काम} \hi{आसान} \hi{होता।}
\end{quote}

\begin{quote}
\emph{agar koī caśmdīd hotā to pulis kā kām āsān hotā.}
\end{quote}

\begin{quote}
If there had been an eyewitness, the police's job would have been easy.
\end{quote}

Masculine or feminine, by the person meant.

\begin{cousinweb}
\textbf{\hi{चश्मदीद}} (\emph{caśmdīd}) is usually traced to Persian \emph{čašm-dīd}, built from \emph{čašm}, \textquotedblleft{}eye\textquotedblright{}, the word behind \textbf{\hi{चश्मा}}, and \emph{dīd}, \textquotedblleft{}seen\textquotedblright{}: a person who saw it with their own eyes.
\end{cousinweb}

\begin{grammarlens}[title={What you've built}]
One more word for telling what happened.
\end{grammarlens}

\subsection*{Your turn}
\begin{itemize}
\raggedright
  \item \emph{Say it:} \emph{caśmdīd}
  \item \emph{Say it:} \emph{caśmdīd}, once more
  \item \emph{Say it:} \emph{agar koī caśmdīd hotā to pulis kā kām āsān hotā.} — If there had been an eyewitness, the police's job would have been easy.
  \item \emph{Recall:} say the Hindi for a police inspector, then the Hindi for a lock-up, then say \emph{caśmdīd} again
\end{itemize}

\begin{tcolorbox}[breakable,colback=teal!4,colframe=teal!35!black,title={Before you move on}]
What does \hi{चश्मदीद} mean? (\textquotedblleft{}An eyewitness\textquotedblright{}.) Say it once more.
\end{tcolorbox}
\section[\texorpdfstring{\hi{बयान} (bayān)}{बयान (bayān)}]{\hi{बयान} (bayān) — a statement}

\label{lesson:HI-C291-bayan}



\begin{quote}
Before the new one: say the Hindi for a lock-up, then the Hindi for an eyewitness.
\end{quote}

\subsection*{You'll want to know: \hi{बयान}}
\textbf{\hi{बयान}} — \emph{bayān} — \textquotedblleft{}a statement\textquotedblright{}.

\begin{quote}
\hi{गवाह} \hi{ने} \hi{थाने} \hi{में} \hi{अपना} \hi{बयान} \hi{दिया।}
\end{quote}

\begin{quote}
\emph{gavāh ne thāne meṁ apnā bayān diyā.}
\end{quote}

\begin{quote}
The witness gave a statement at the police station.
\end{quote}

Masculine.

\begin{cousinweb}
\textbf{\hi{बयान}} (\emph{bayān}) is usually traced, through Persian, to Arabic \emph{bayān}, \textquotedblleft{}an explanation, a making clear\textquotedblright{}. \textbf{\hi{बयान} \hi{देना}} is to give a statement.
\end{cousinweb}

\begin{grammarlens}[title={What you've built}]
One more word for telling what happened.
\end{grammarlens}

\subsection*{Your turn}
\begin{itemize}
\raggedright
  \item \emph{Say it:} \emph{bayān}
  \item \emph{Say it:} \emph{bayān}, once more
  \item \emph{Say it:} \emph{gavāh ne thāne meṁ apnā bayān diyā.} — The witness gave a statement at the police station.
  \item \emph{Recall:} say the Hindi for a lock-up, then the Hindi for an eyewitness, then say \emph{bayān} again
\end{itemize}

\begin{tcolorbox}[breakable,colback=teal!4,colframe=teal!35!black,title={Before you move on}]
What does \hi{बयान} mean? (\textquotedblleft{}A statement\textquotedblright{}.) Say it once more.
\end{tcolorbox}
\section[\texorpdfstring{\hi{पहचानना} (pahcānnā)}{पहचानना (pahcānnā)}]{\hi{पहचानना} (pahcānnā) — to recognize}

\label{lesson:HI-C291-pahcanna}



\begin{quote}
Before the new one: say the Hindi for an eyewitness, then the Hindi for a statement.
\end{quote}

\subsection*{You'll want to know: \hi{पहचानना}}
\textbf{\hi{पहचानना}} — \emph{pahcānnā} — \textquotedblleft{}to recognize\textquotedblright{}.

\begin{quote}
\hi{मीरा} \hi{ने} \hi{चोर} \hi{को} \hi{पहचान} \hi{लिया।}
\end{quote}

\begin{quote}
\emph{Mīrā ne cor ko pahcān liyā.}
\end{quote}

\begin{quote}
Mira recognized the thief.
\end{quote}

Past, with \hi{ने} (he or she …): \hi{उसने} \hi{चोर} \hi{को} \hi{पहचाना}, \emph{usne cor ko pahcānā}.

\begin{cousinweb}
\textbf{\hi{पहचानना}} (\emph{pahcānnā}) goes back to Sanskrit \emph{pratyabhijānāti}, \textquotedblleft{}recognizes\textquotedblright{} (CDIAL 8637). The noun \textbf{\hi{पहचान}}, identity or recognition, belongs with it; the verb is what you do when a face is one you know.
\end{cousinweb}

\begin{grammarlens}[title={What you've built}]
One more word for telling what happened.
\end{grammarlens}

\subsection*{Your turn}
\begin{itemize}
\raggedright
  \item \emph{Say it:} \emph{pahcānnā}
  \item \emph{Say it:} \emph{pahcānnā}, once more
  \item \emph{Say it:} \emph{Mīrā ne cor ko pahcān liyā.} — Mira recognized the thief.
  \item \emph{Recall:} say the Hindi for an eyewitness, then the Hindi for a statement, then say \emph{pahcānnā} again
\end{itemize}

\begin{tcolorbox}[breakable,colback=teal!4,colframe=teal!35!black,title={Before you move on}]
What does \hi{पहचानना} mean? (\textquotedblleft{}To recognize\textquotedblright{}.) Say it once more.
\end{tcolorbox}
\section[\texorpdfstring{\hi{शक} (śak)}{शक (śak)}]{\hi{शक} (śak) — suspicion; doubt}

\label{lesson:HI-C291-sak}



\begin{quote}
Before the new one: say the Hindi for a statement, then the Hindi for to recognize.
\end{quote}

\subsection*{You'll want to know: \hi{शक}}
\textbf{\hi{शक}} — \emph{śak} — \textquotedblleft{}suspicion; doubt\textquotedblright{}.

\begin{quote}
\hi{मुझे} \hi{उस} \hi{आदमी} \hi{पर} \hi{शक} \hi{है।}
\end{quote}

\begin{quote}
\emph{mujhe us ādmī par śak hai.}
\end{quote}

\begin{quote}
I suspect that man.
\end{quote}

Masculine.

\begin{cousinweb}
\textbf{\hi{शक}} (\emph{śak}) is usually traced, through Persian, to Arabic \emph{shakk}, \textquotedblleft{}doubt\textquotedblright{}. The person you suspect takes \textbf{\hi{पर}}: \textbf{\hi{मुझे} \hi{उस} \hi{पर} \hi{शक} \hi{है}}.
\end{cousinweb}

\begin{grammarlens}[title={What you've built}]
One more word for telling what happened.
\end{grammarlens}

\subsection*{Your turn}
\begin{itemize}
\raggedright
  \item \emph{Say it:} \emph{śak}
  \item \emph{Say it:} \emph{śak}, once more
  \item \emph{Say it:} \emph{mujhe us ādmī par śak hai.} — I suspect that man.
  \item \emph{Recall:} say the Hindi for a statement, then the Hindi for to recognize, then say \emph{śak} again
\end{itemize}

\begin{tcolorbox}[breakable,colback=teal!4,colframe=teal!35!black,title={Before you move on}]
What does \hi{शक} mean? (\textquotedblleft{}Suspicion; doubt\textquotedblright{}.) Say it once more.
\end{tcolorbox}
\section[\texorpdfstring{\hi{घूरना} (ghūrnā)}{घूरना (ghūrnā)}]{\hi{घूरना} (ghūrnā) — to stare}

\label{lesson:HI-C291-ghurna}



\begin{quote}
Before the new one: say the Hindi for to recognize, then the Hindi for suspicion.
\end{quote}

\subsection*{You'll want to know: \hi{घूरना}}
\textbf{\hi{घूरना}} — \emph{ghūrnā} — \textquotedblleft{}to stare\textquotedblright{}.

\begin{quote}
\hi{वह} \hi{आदमी} \hi{मुझे} \hi{घूर} \hi{रहा} \hi{था।}
\end{quote}

\begin{quote}
\emph{vah ādmī mujhe ghūr rahā thā.}
\end{quote}

\begin{quote}
That man was staring at me.
\end{quote}

Past, with \hi{ने} (he or she …): \hi{उसने} \hi{मुझे} \hi{घूरा}, \emph{usne mujhe ghūrā}.

\begin{cousinweb}
\textbf{\hi{घूरना}} (\emph{ghūrnā}) comes from \emph{*ghūra}, a frown, a form CDIAL reconstructs from its descendants (CDIAL 4496). It is a hard, long look, never a friendly one.
\end{cousinweb}

\begin{grammarlens}[title={What you've built}]
One more word for telling what happened.
\end{grammarlens}

\subsection*{Your turn}
\begin{itemize}
\raggedright
  \item \emph{Say it:} \emph{ghūrnā}
  \item \emph{Say it:} \emph{ghūrnā}, once more
  \item \emph{Say it:} \emph{vah ādmī mujhe ghūr rahā thā.} — That man was staring at me.
  \item \emph{Recall:} say the Hindi for to recognize, then the Hindi for suspicion, then say \emph{ghūrnā} again
\end{itemize}

\begin{tcolorbox}[breakable,colback=teal!4,colframe=teal!35!black,title={Before you move on}]
What does \hi{घूरना} mean? (\textquotedblleft{}To stare\textquotedblright{}.) Say it once more.
\end{tcolorbox}
